\chapter{Revisão de Literatura}
\label{cap:rsl}

Este capítulo apresenta a Revisão Sistemática da Literatura (RSL) conduzida para
fundamentar a modernização do módulo GPSR no ns-3. Diferentemente de uma revisão
narrativa, a RSL segue um protocolo definido previamente à busca, o que torna o processo de
seleção e análise dos estudos transparente e reproduzível \cite{kitchenham2007guidelines}.
Como referência de escopo, toma-se o estudo de \citeonline{campanile2020computer}, que
mapeou o uso do ns-3 em publicações de 2009 a 2019; a presente revisão amplia esse
levantamento, concentrando-se na extensibilidade do simulador no domínio das redes
veiculares.

A revisão é guiada pelas seguintes questões de pesquisa (\textit{Research Questions}, RQ):

\begin{itemize}
  \item \textbf{RQ1:} Que extensões e módulos voltados a redes veiculares já foram propostos
  ou implementados no ns-3?
  \item \textbf{RQ2:} Que protocolos de roteamento para VANETs foram implementados ou
  avaliados no ns-3, e em que medida o roteamento geográfico (GPSR e variantes) está
  representado?
  \item \textbf{RQ3:} Que lacunas funcionais, limitações de compatibilidade ou dificuldades
  de manutenção de módulos são reportadas pela literatura?
  \item \textbf{RQ4:} Que metodologias de validação e métricas de desempenho são empregadas
  para avaliar módulos e protocolos no ns-3?
\end{itemize}

\section{Protocolo da Revisão}
\label{sec:protocolo-rsl}

O protocolo compreende as bases pesquisadas, as \textit{strings} de busca, o período
considerado e o registro da execução, descritos a seguir.

\subsection{Bases Pesquisadas}
\label{sec:bases}

A busca é realizada nas bases IEEE Xplore, ACM Digital Library, Scopus e Web of Science,
por concentrarem a produção científica da área de redes de computadores. De forma
complementar, utiliza-se o Google Scholar para identificar trabalhos não indexados nas
bases anteriores e a biblioteca digital da Sociedade Brasileira de Computação (SBC
OpenLib), a fim de incluir a produção nacional, como os anais do Simpósio Brasileiro de Redes
de Computadores e Sistemas Distribuídos (SBRC).

\subsection{\textit{Strings} de Busca}
\label{sec:strings}

A \textit{string} base (S1) combina três grupos de termos: o simulador, o domínio de aplicação
e o tipo de contribuição. Uma segunda \textit{string} (S2) é dirigida ao roteamento geográfico,
e uma terceira (S3) é utilizada nas bases em português:

\begin{quote}
\small
\textbf{S1:} \texttt{("ns-3" OR "ns3" OR "network simulator 3") AND ("VANET" OR "vehicular
ad hoc" OR "vehicular network*") AND ("routing" OR "module" OR "extension" OR
"implementation")}

\vspace{0.2cm}
\textbf{S2:} \texttt{("ns-3" OR "ns3") AND ("GPSR" OR "geographic routing" OR
"position-based routing")}

\vspace{0.2cm}
\textbf{S3:} \texttt{("ns-3" OR "ns3") AND ("VANET" OR "redes veiculares") AND
("roteamento" OR "módulo" OR "extensão")}
\end{quote}

As \textit{strings} são adaptadas à sintaxe de cada base (campos de busca, operadores e
truncamento) e aplicadas aos campos de título, resumo e palavras-chave. Como verificação de
sensibilidade, confere-se se a busca recupera trabalhos já conhecidos previamente pelo autor
(conjunto-semente), ajustando-se os termos caso algum deles não seja retornado.

\subsection{Período Considerado}
\label{sec:periodo}

São considerados os estudos publicados entre 2009 e 2026, de modo a cobrir a janela
analisada por \citeonline{campanile2020computer} e o intervalo posterior a ela.

\subsection{Data da Pesquisa}
\label{sec:data-pesquisa}

As buscas são executadas em um único período contínuo, para reduzir variações nos
resultados retornados pelas bases ao longo do tempo. A data de execução e o número de
resultados retornados em cada base são registrados na Tabela~\ref{tab:busca}.

\begin{table}[ht]
\centering
\caption{Registro da execução da busca por base.}
\label{tab:busca}
\small
\begin{tabular}{|l|c|c|}
\hline
\textbf{Base} & \textbf{Data da busca} & \textbf{Resultados retornados} \\ \hline
IEEE Xplore & -- & -- \\ \hline
ACM Digital Library & -- & -- \\ \hline
Scopus & -- & -- \\ \hline
Web of Science & -- & -- \\ \hline
Google Scholar & -- & -- \\ \hline
SBC OpenLib & -- & -- \\ \hline
\end{tabular}

\par\vspace{0.1cm}{\small Fonte: elaborado pelo autor.}
\end{table}

\section{Critérios de Inclusão e Exclusão}
\label{sec:criterios}

Consideram-se apenas estudos completos, publicados em periódicos ou anais de eventos, em
português ou inglês, dentro do período definido na Seção~\ref{sec:periodo}.

\subsection{Critérios de Inclusão}
\label{sec:inclusao}

Um estudo é incluído quando atende ao critério CI1 e a pelo menos um dos critérios CI2 a CI5:

\begin{itemize}
  \item \textbf{CI1:} envolve o ns-3 como ferramenta de simulação;
  \item \textbf{CI2:} trata de redes veiculares (VANETs);
  \item \textbf{CI3:} aborda o protocolo GPSR ou o roteamento geográfico;
  \item \textbf{CI4:} discute a extensibilidade do ns-3;
  \item \textbf{CI5:} descreve a implementação de módulos no ns-3.
\end{itemize}

\subsection{Critérios de Exclusão}
\label{sec:exclusao}

Um estudo é excluído quando atende a qualquer um dos critérios abaixo:

\begin{itemize}
  \item \textbf{CE1:} estudo duplicado;
  \item \textbf{CE2:} texto sem acesso integral;
  \item \textbf{CE3:} sem relação com o tema da pesquisa;
  \item \textbf{CE4:} cita o ns-3 sem utilizá-lo como ferramenta.
\end{itemize}

\section{Seleção dos Estudos e Extração de Dados}
\label{sec:selecao}

A seleção é conduzida em quatro etapas sucessivas: (i) execução das \textit{strings} em cada
base e remoção de duplicatas; (ii) triagem por título e resumo, aplicando os critérios da
Seção~\ref{sec:criterios}; (iii) leitura completa dos estudos remanescentes, com nova
aplicação dos critérios; e (iv) extração dos dados dos estudos incluídos. O número de
estudos em cada etapa é registrado na Tabela~\ref{tab:fluxo}.

\begin{table}[ht]
\centering
\caption{Fluxo de seleção dos estudos.}
\label{tab:fluxo}
\small
\begin{tabular}{|l|c|}
\hline
\textbf{Etapa} & \textbf{Estudos} \\ \hline
Registros identificados nas bases & -- \\ \hline
Após remoção de duplicatas & -- \\ \hline
Após triagem por título e resumo & -- \\ \hline
Incluídos após leitura completa & -- \\ \hline
\end{tabular}

\par\vspace{0.1cm}{\small Fonte: elaborado pelo autor.}
\end{table}

Para cada estudo incluído, são extraídos: autores e ano; versão do ns-3 utilizada; protocolo ou
módulo estudado; se houve criação ou modificação de módulo; modelo de mobilidade
empregado; métricas de desempenho avaliadas; e limitações ou problemas de compatibilidade
relatados. Esses campos alimentam a resposta às questões RQ1 a RQ4.

\section{Análise dos Estudos Selecionados}
\label{sec:analise-rsl}

% A ser redigida após a execução do protocolo: síntese por questão de pesquisa (RQ1 a RQ4).

\section{Lacunas Identificadas}
\label{sec:lacunas}

% A ser redigida após a análise dos estudos selecionados.